\chapter{Discussion}
\label{chap:discussion}

\section{Overview}
\label{sec:discussion-overview}

Chapter~\ref{chap:results} reported the quantitative results of the four-pipeline comparison. This chapter interprets what those results mean. We discuss why prompt improvements alone are not enough to close the quality gap, how knowledge graph data can be used more effectively, how the approach may transfer to other domains, and where the pipeline falls short. Chapter~\ref{chap:conclusion} then answers each research question and outlines future work.


\section{Why Prompt Changes Are Not Enough}
\label{sec:discuss-prompts}

The no-ontology pipeline and the two-step baseline both retrieve from the same knowledge graph. The difference is that the no-ontology variant uses improved prompts for refinement and generation, organised around the three-agent design with structured output (Topics \& Information, Validated References, Caveats \& Gaps). This gives it a 1.37-point advantage over the two-step baseline (6.88 vs 5.51), which is a large improvement. Prompt changes clearly help. They fix the structural and directness problems that the two-step baseline's raw context injection causes. But the no-ontology pipeline is only marginally better than zero-shot on core quality metrics: relevance is 7.56 vs 7.57 and completeness is 6.86 vs 6.91. The improved prompts recover what bad context processing destroyed, but they do not push quality beyond what the model can already produce on its own.

The ontology adds something that prompt tuning cannot: a query-specific semantic structure that tells the refinement agent what to look for in the documents, and tells the generation agent what dimensions to cover in the answer. Without the ontology, the refinement agent has to decide on its own which parts of the context are relevant to the query. With the ontology, it has a checklist: does this document address the biophysical variable? The sensor? The spatial scope? This is what produces the additional +0.60 on top of the already-improved pipeline. The largest gains are on helpfulness (+0.86) and directness (+0.81), which are the criteria where knowing what matters most makes the biggest difference. As~\citet{liu2024lost} showed, even well-prompted models lose information in long contexts. The ontology acts as an anchor that prevents this by marking what is important before the context is even constructed. In short, prompts change how the model processes context; the ontology changes what context the model gets to process.


\section{Making Use of Knowledge Graph Data}
\label{sec:discuss-kg-data}

The main lesson from the results is that context quality matters more than context quantity. The two-step baseline provides shorter context (about 2{,}948 characters on average) but in unstructured form, and the answers are worse than zero-shot. The three-agent pipeline provides longer context (about 9{,}400--9{,}600 characters) but in structured form, and the answers are better. The knowledge graph built by~\citet{elbaff2025earthobservationrag} is a strong retrieval source. The AQL queries return relevant documents. The problem was never what gets retrieved, but what happens after retrieval: how the context is filtered, organised, and presented to the generation model. This is a general insight for any RAG system built on a knowledge graph. The graph gives you relevant data, but the pipeline has to structure that data around the query's specific needs before passing it to the generator.

One concrete data limitation is that the knowledge graph currently indexes only document abstracts. Abstracts summarise findings in general terms and rarely include exact measurements or detailed statistics. This is the main reason quantitative precision is low across all four pipelines (4.01--5.93). The ontology helps to some extent: centrality scores flag when quantitative detail should exist for a given topic, prompting the refinement agent to look for it. But if the numbers are not in the abstracts, the pipeline cannot surface them. Indexing full-text sections, especially results tables and method descriptions, would give the pipeline access to the specific numbers it currently lacks. This is a data change rather than a pipeline change; the architecture is already ready for richer input.


\section{Applicability Beyond Earth Observation}
\label{sec:discuss-generalisability}

The ontology agent is dynamic: it generates a new ontology for each query rather than relying on a pre-built, fixed ontology. The five semantic dimensions we use for EO (biophysical variable, sensor/method, spatial scope, temporal scope, application domain) are one way to instantiate this. Any domain with structured knowledge has similar dimensions. In medical QA these could be patient population, biomarker, intervention, temporal window, and clinical setting. In legal QA: jurisdiction, legal instrument, temporal applicability, and subject matter. In engineering: material, measurement method, environmental conditions, scale, and application domain. As long as the ontology agent's prompt is adapted to define the right dimensions for the new domain, and the retrieval backend provides good source data, the same three-agent architecture should work. The refinement and generation agents already operate on the ontology's output rather than on hard-coded EO knowledge.

Two things are needed for domain transfer: a quality retrieval source with good domain coverage, and a set of semantic dimensions that capture how questions in that domain are structured. We were fortunate to have a well-curated knowledge graph with over 2 million documents~\citep{elbaff2025earthobservationrag}. The pipeline's performance depends on this. A weaker retrieval source would likely reduce the ontology's benefit, because the refinement agent would have less material to work with. Adapting the ontology dimensions is the main effort, and it requires domain understanding. Someone needs to define what the key semantic axes of the domain are. Once that is done, the ontology agent prompt is a relatively small change and the rest of the pipeline stays the same. We should be clear that we tested on one domain, one knowledge graph, and one model. The argument for generalisability is structural, not empirical, and validation in other domains remains future work.


\section{Limitations}
\label{sec:limitations}

\paragraph{Quantitative precision.} The best pipeline still scores only 5.93/10 on quantitative precision. The root cause is that the knowledge graph indexes only abstracts, which rarely contain specific numbers. For researchers who need exact measurements or statistical values, this is a real limitation.

\paragraph{Vague and incomplete queries.} All 70~evaluation questions are well-formed and specific. We did not test what happens when a user asks something vague like ``tell me about forests'' or provides an incomplete prompt. The ontology agent assumes a clear query to decompose. With a vague query, it may produce an overly broad or unfocused ontology, which would reduce its filtering value. This is an untested failure mode.

\paragraph{Single model, single knowledge graph, single judge.} All pipelines use Mistral Small. The results may not generalise to other model families. The knowledge graph comes from a single source. The judge is one model (Claude Sonnet~4.6). These are standard constraints for a bachelor thesis, but they limit how far the findings can be generalised.

\paragraph{Evaluation framework.} The LLM-as-a-judge approach inherits the judge model's biases, and context-blind judging means we cannot verify whether claims are actually grounded in retrieved documents. The moderate inter-criterion correlation ($r = 0.65$) also means the eight criteria are not fully independent. Chapter~\ref{chap:evaluation} discusses these issues in detail.

\paragraph{No user study.} All evaluation is automated. We do not know whether EO researchers would agree with the judge model's scores or find the answers useful in practice.
